% ==========================================================================
%  Section 8 -- Cylinder uniqueness and identification.  Owner: W08.
%  Frozen labels: sec:identification, thm:cylinder-uniqueness,
%  lem:subsequence-marginals.
%
%  Two statements, kept apart on purpose.  thm:cylinder-uniqueness is about
%  the carrier alone: two Borel PROBABILITY measures that agree on every
%  finite cylinder set coincide, and it is the probability hypothesis that
%  makes the pi-lambda step valid.  lem:subsequence-marginals is about the
%  sequence alone: a limit point extracted from (mu_N) inherits the target
%  marginals of thm:finite-family.  Neither is a uniqueness statement about
%  a filter limit; the characterisation lives in sec:main and is proved
%  there from these two results.
% ==========================================================================
\section{Cylinder uniqueness and identification}\label{sec:identification}

This section supplies the two statements that turn a tight family into one
identified limit, and it keeps them separate because they answer different
questions. \Cref{thm:cylinder-uniqueness} is a statement about the carrier
alone: two Borel probability measures there which agree on every finite cylinder
set are equal. \Cref{lem:subsequence-marginals} is a statement about the
sequence alone: a weak limit point extracted from $(\mu_N)$ inherits the finite
marginals that \cref{thm:finite-family} establishes for the sequence. The first
says that the target marginal family admits at most one measure; the second says
that every limit point is a measure admitted by it. Together they say that the
limit points of the sequence are one point, which is the conclusion
\cref{sec:main} draws. Neither statement identifies the limit with an operator,
and neither is used as if it did.

\subsection{Determination by cylinders}

Throughout, a \emph{cylinder set} is a set of the form
$\Pi_F^{-1}(B)\subseteq\mathcal C$, where $F$ is a finite tuple of modes,
$\Pi_F$ is the finite-mode coordinate map of \cref{def:cylinder-map} and
$B$ is a Borel subset of the finite-dimensional target of $\Pi_F$.

\begin{theorem}[Cylinder-to-Borel uniqueness]\label{thm:cylinder-uniqueness}
Let $\mu$ and $\nu$ be Borel \emph{probability} measures on the carrier
$\mathcal C$. If $\mu(A)=\nu(A)$ for every cylinder set $A$, then $\mu=\nu$.
Equivalently: if all finite-coordinate pushforward laws of $\mu$ and $\nu$
coincide, then $\mu=\nu$.
\end{theorem}

\begin{proof}
\emph{(i) The cylinders form a $\pi$-system of Borel sets.} Let
$A=\Pi_F^{-1}(B)$ and $A'=\Pi_{F'}^{-1}(B')$ be cylinder sets. The map
$\Pi_{F\cup F'}$ of the concatenated tuple has $\Pi_F$ and $\Pi_{F'}$ among its
coordinate projections, so $A\cap A'$ is the preimage under $\Pi_{F\cup F'}$ of
a Borel set, hence a cylinder set. Each cylinder set is Borel, because $\Pi_F$ is
continuous (\cref{lem:coordinate-embedding}); so this family is a $\pi$-system
contained in $\mathcal B(\mathcal C)$.

\emph{(ii) It generates the Borel sigma algebra.} Two inclusions.
\emph{$\subseteq$:} each $\Pi_F$ is continuous (\cref{lem:coordinate-embedding}), so
$\sigma(\Pi_F)\subseteq\mathcal B(\mathcal C)$.
\emph{$\supseteq$:} the tuple $((k,\mathrm c))$ is admissible, so
$\Re c_k=\Pi_{((k,\mathrm c))}$ is $\sigma(\Pi_F)$-measurable; likewise
$-\Im c_k=\Pi_{((k,\mathrm s))}$, again for every $k$, the zero mode included.
Hence each coordinate functional $\pi_k=\Re c_k+i\Im c_k$ is measurable for
$\sigma(\Pi_F)$, so $\sigma(\pi_k:k\in\Z^2)\subseteq\sigma(\Pi_F)$, and
$\sigma(\pi_k)=\mathcal B(\mathcal C)$ by \cref{thm:cylinder-borel}; separability
enters through that last step, via \cref{thm:separable}.
The second inclusion is the only place where the definition of
\cref{def:cylinder-map} is used, and it is not decorative. Take the restricted
family that keeps the cosine label at every mode and drops the sine label at
$k=0$. The two Dirac measures $\delta_{0}$ and $\delta_{ie_0}$ on
$\mathcal C$ --- where $e_0$ is the unit vector at the zero mode --- agree on
every slot of that family ($\Re c_0=0$ at both points, and every $c_k=0$ for
$k\ne0$) and differ in the dropped slot ($-\Im c_0=0$ against $-1$). So
\cref{thm:cylinder-uniqueness} would fail. The same construction applies to the
omission of any single real coordinate: the pair of Dirac masses at two points
that differ only in that coordinate separates the restricted family.

\emph{(iii) The Dynkin step.} Fix $\mu$ and put
$\mathcal D:=\{A\in\mathcal B(\mathcal C):\mu(A)=\nu(A)\}$. It contains every
cylinder set by hypothesis. It contains $\mathcal C$ itself, because
$\mu(\mathcal C)=\nu(\mathcal C)=1$. It is a $\lambda$-system: if $A\subseteq A'$
both lie in $\mathcal D$ then
$\mu(A'\setminus A)=\mu(A')-\mu(A)=\nu(A')-\nu(A)=\nu(A'\setminus A)$, and if
$A_1,A_2,\dots\in\mathcal D$ are pairwise disjoint then countable additivity of
both measures gives
$\mu(\bigcup_jA_j)=\sum_j\mu(A_j)=\sum_j\nu(A_j)=\nu(\bigcup_jA_j)$. By Dynkin's
$\pi$--$\lambda$ theorem a $\lambda$-system containing a $\pi$-system contains
the sigma algebra that $\pi$-system generates, so by (i) and (ii)
$\mathcal D=\mathcal B(\mathcal C)$, that is $\mu=\nu$.

\emph{(iv) Where the word probability is used, and what the conclusion buys.}
Steps (ii) and (iii) hold for any two finite measures of equal total mass; it is
the hypothesis that both are probability measures which makes that total mass
equal, and the hypothesis cannot be dropped. For two infinite measures agreeing
on every cylinder set, the difference in Step (iii) is undefined at the first
complement, and equality of total mass is exactly what is missing. From
$\mu=\nu$ we obtain equality of integrals of every bounded Borel function, and
every bounded continuous function $\mathcal C\to\R$ is Borel because it is
continuous. The sigma algebra equality (ii) is what makes ``agrees on every
cylinder'' and ``has the same law'' the same requirement.
\end{proof}

\subsection{A subsequential limit inherits the target marginals}

The second statement is the only place where the marginal theorem of
\cref{sec:marginals} meets the extraction theorem of \cref{sec:tightness}. It
says nothing about how many limit points there are.

\begin{lemma}[A subsequential limit inherits the target marginals]\label{lem:subsequence-marginals}
Let $L>0$, let $m>0$, let $(\mu_N)$ be the sequence \eqref{eq:finite-source-law}
of finite-source Fourier pushforwards, let $(N_j)$ be a strictly increasing
sequence of indices and suppose that $\mu_{N_j}\Rightarrow\mu$ weakly on
$\mathcal C$. Then for every finite tuple $F$ of modes the pushforward
$(\Pi_F)_*\mu$ is the resolvent Gaussian law $\Lambda_F$ that
\cref{thm:finite-family} records as the limit of the finite-mode laws
$(\Pi_F)_*\mu_N$.
\end{lemma}

\begin{proof}
Fix $F$, and abbreviate $\mu_N^F:=(\Pi_F)_*\mu_N$ and
$\mu^F:=(\Pi_F)_*\mu$. The map $\Pi_F$ is continuous
(\cref{lem:coordinate-embedding}), and a continuous image of a weakly convergent
sequence of Borel probability measures converges weakly to the image of the
limit: for every bounded continuous $\psi$ on the finite-dimensional target, the
composition $\psi\circ\Pi_F$ is bounded and continuous on $\mathcal C$, so
\[
  \int\psi\,\dd\mu_{N_j}^F
  =\int(\psi\circ\Pi_F)\,\dd\mu_{N_j}
  \longrightarrow\int(\psi\circ\Pi_F)\,\dd\mu
  =\int\psi\,\dd\mu^F .
\]
Hence $\mu_{N_j}^F\Rightarrow\mu^F$. For the second limit, note what
$\mu_N^F$ is: the stored coordinate at a mode is $\sqrt{w_3}$ times the Fourier
coordinate, so its law is the image of the law of the real cosine and sine
coefficients of that mode under the fixed reassembly map which multiplies by
$\sqrt{w_3}$ (\cref{def:cylinder-map}). By \cref{thm:finite-family} the laws of
those real coefficients converge weakly, for every finite family of test
functions, to the resolvent Gaussian law; a fixed continuous reassembly carries
that convergence to $\Lambda_F$, and a subsequence of a convergent sequence has
the same limit. Hence $\mu_{N_j}^F\Rightarrow\Lambda_F$. The target of $\Pi_F$
is a metric space, so its Borel probability measures with the weak topology form
a Hausdorff space, in which a sequence has at most one limit; therefore
$\mu^F=\Lambda_F$.
\end{proof}

These two results are the last probabilistic steps of Part~I. What stands
between them and the theorem is purely topological: at this point one knows that
every subsequence of $(\mu_N)$ has a further subsequence converging to one and
the same measure, and the passage from that statement to convergence of the
whole sequence mentions no measure and no probability at all. \Cref{sec:main}
makes that passage, as Step~5 of the proof of \cref{thm:main}, and names the
external fact it uses (\cref{ext:subseq}).
